% !TeX root = ../../Book3.tex
% 纲要标题：### F.4 长度、重数与局部消解的计算

\section{长度、重数与局部消解的计算}
\label{b3:sec:F04}

局部长度与交数的计算必须证明高阶项不再贡献，否则得到的只是一个截断值。下面把理想幂次包含、Jacobian 商和自由消解作为可核验的停止与结构证据。

% 纲要细目（与计划纲要同步）：
% * 有限长度、Hilbert–Samuel 函数及重数的计算证书；记录系数域、工作环、项序与截断精度
% * 模的标准基、合冲模与局部自由消解仅作简要介绍；第21章的全局 Schreyer 方法在正文独立建立
% * 与附录 D 的幂级数结构合用；完整局部算法、复杂度和软件实验可独立编成计算代数或奇点计算专著
%
% #### F.4.1 有限长度、Hilbert–Samuel 函数与重数的计算
% #### F.4.2 模的标准基、合冲模与局部自由消解
% #### F.4.3 幂级数结构与局部算法专题导读


\subsection{有限长度、Hilbert–Samuel 函数与重数的计算}
\label{b3:subsec:F04:01}

有限长度的计算要有一个停止依据。若只算到 \(N\) 阶，没有说明更高阶为何不再贡献，就只能得到截断商的长度。

\begin{proposition}\label{b3:prop:F-truncation-length}
设 \(k\) 为域，\(R=k[x_1,\ldots,x_n]_{(x)}\) 或 \(k[[x_1,\ldots,x_n]]\)，其中 \((x)=(x_1,\ldots,x_n)\)，\(I\subset R\) 为真理想，\(N\ge1\)。若 \((x)^N\subset I\)，则
\[
R/I\cong R/(I+(x)^N),
\qquad
\ell_R(R/I)=\dim_k(R/I).
\]
因此，次数小于 \(N\) 的单项式及其模 \(I+(x)^N\) 的全部线性关系足以确定完整长度。若只知道生成元模 \((x)^N\) 的值而没有上述包含关系，则同一计算一般只确定右边的截断商。
\end{proposition}
\begin{proof}
包含关系使两个商相同。有限维局部 \(k\)-代数的唯一简单模是 \(k\)，故一条组成列的每个因子在 \(k\) 上为一维，长度等于向量空间维数。模 \((x)^N\) 后只剩有限多个单项式；理想的像由所有生成元与单项式的乘积、再截去次数至少为 \(N\) 的项张成，所以精确行消元给出全部线性关系。
\end{proof}

例如\cref{b3:exm:F-two-equations}已经证明 \(\mathfrak m^4=0\)，故四阶截断给出精确长度 \(5\)。另一个常用证书是给出各次数 \(N\) 单项式在 \(I\) 中的表示；有限个这样的表示便证明 \((x)^N\subset I\)。如果先证明商模的 \(\mathfrak m^N\) 与 \(\mathfrak m^{N+1}\) 相同，Noether 性与 Nakayama 引理也可推出 \(\mathfrak m^N=0\)。但截断数据中“没有看到新项”本身不是这个模等式的证明。

\ProseParagraphBreak
正维情形要计算分次层的持续增长。\cref{b3:thm:hilbert-samuel}保证，若 \(A=R/I\) 的维数为 \(d\)，那么
\[
\ell_A(A/\mathfrak m^{q+1})
=\frac{e(A)}{d!}q^d+\text{低次项}\qquad(q\gg0).
\]
若已用标准基得到切锥理想，就可按单项式或齐次理想的 Hilbert 级数计算它；\cref{b3:prop:B-monomial-hilbert}给出单项式情形的有限公式。这里必须确认得到的是全部初始理想，否则错误的分次商可能连维数都不对。

\begin{example}\label{b3:exm:F-jacobian-quotients}
设 \(k\) 特征为零，\(R=k[x,y]_{(x,y)}\)，\(f=y^2-x^3\)。其偏导数理想为
\[
J_f=(\partial f/\partial x,\partial f/\partial y)=(x^2,y).
\]
因此 \(R/J_f\) 以 \(1,x\) 为基，长度为 \(2\)。加入 \(f\) 不改变理想，所以 \(R/(f,J_f)\) 也长为 \(2\)。在孤立超曲面奇点的语言中，这两种有限长度分别称为\term{Milnor 数}[Milnor number]和\term{Tjurina 数}[Tjurina number]；此处仅计算这些代数商，不引入复解析拓扑解释。\chapref{b3:ch:14}的微分展示说明 Jacobian 矩阵为何出现，而本附录说明如何验证由它得到的有限长度。特征条件不可省：若特征为 \(3\)，则 \(J_f=(y)\)，第一个商已不具有有限长度。
\end{example}

\subsection{模的标准基、合冲模与局部自由消解}
\label{b3:subsec:F04:02}

模的局部标准基与理想的差别，在于首项还须记录所在的基分量。设 \(R=k[x]_{(x)}\)，在 \(R^r\) 的单项式 \(x^\alpha e_j\) 上选定与乘法相容的局部模次序。对 \(N\subset R^r\)，用全部非零元素的首项生成 \(k[x]^r\) 的单项式子模；有限生成组与它具有相同首项子模时，称为模标准基。单位乘子和弱正规形同样可逐分量定义。

\ProseParagraphBreak
计算合冲时，仍须先记清楚原生成矩阵。若 \(G=FU,\ F=GV\) 是局部环中的两向变换，且 \(T\) 生成 \(\ker G\)，则与\secref{b3:sec:B04}完全相同的恒等式给出
\[
\ker F=\operatorname{im}[\,UT\ \ I-UV\,].
\]
这条矩阵等式不依赖全局序；但“怎样求得生成全部 \(\ker G\) 的 \(T\)”依赖正确的局部标准基方法，不能把全局 Schreyer 算法的终止证明原样搬来。

\begin{example}\label{b3:exm:F-local-resolution}
取 \(R=k[x,y]_{(x,y)}\)。理想 \((x,y)\) 的全部关系由 \((-y,x)^{\mathsf T}\) 生成，因此
\[
0\longrightarrow R
\xrightarrow{\binom{-y}{x}}R^2
\xrightarrow{(x\ \ y)}R
\longrightarrow k\longrightarrow0
\]
是一个自由消解。确实，若 \(ax+by=0\)，则模 \((x)\) 后，整环 \(R/(x)\cong k[y]_{(y)}\) 中有 \(\bar b\,y=0\)，故 \(b=xc\)，再约去 \(x\) 得 \(a=-yc\)。左映射单射来自 \(R\) 为整环。所有微分项属于极大理想，故该消解极小，Betti 数为 \(1,2,1\)。
\end{example}

局部极小化时，任何可逆矩阵项都可作主元，而不只非零常数。例如 \(1-x\) 在多项式环中不是可逆项，在原点局部环中却可以消去一个可缩对。消去后必须同时更新左右相邻微分。对于奇异局部商环，极小自由消解还可能无限长：在 \(A=k[\epsilon]/(\epsilon^2)\) 上，剩余域有无穷消解，其每个正次微分都是乘 \(\epsilon\)，因为该映射的核和像都为 \((\epsilon)\)。这与\chapref{b3:ch:21}多项式环上的有限长度结论并不矛盾。

\subsection{幂级数结构与局部算法专题导读}
\label{b3:subsec:F04:03}

本附录中的局部计算由三种有限信息支撑：原生成元之间的精确等式、经过核验的首项理想，以及保证截断已经足够的理想幂包含关系。有了这些信息，就可以逐项核验理想成员、切锥、长度和某些消解；缺少其中一项时，输出可能只是一份候选或有限精度近似。

\ProseParagraphBreak
\appref{b3:app:D}进一步讨论局部环与幂级数环之间的比较、Cohen 结构以及 Weierstrass 方法。结构定理解释哪些完备局部环能写成幂级数环的商；标准基方法则从一组给定方程计算这个商的初始关系。二者的假设各有作用：给出形式表示不等于给出全部系数的有限算法，给出有限截断也不等于决定一般级数是否为零。

\ProseParagraphBreak
若继续研究局部算法，需要系统处理一般局部序和混合序、Mora 弱正规形的终止性、标准基的临界对判据、局部合冲计算与有效界。若转向奇点计算，还需增加分支、等价关系和所用不变量的几何意义。这些问题值得专门学习，本卷只选择能够独立核验的基本构造与实例，不把一个软件返回的整数当作未说明的理论结论。
